\section*{ID6684: Bessel + Aperture → Spectral Cutoff, Link to Spherical Harmonics: D/λ}
\paragraph{Metadata.}
PiID: 3005450580685 | RTEID: RTE:P25955.4564:T5.0675 | UUID: 032efd44-0116-5b23-9abc-f78f76a538d1

\paragraph{Canonical Statement.}
So the **aperture diameter measured in wavelengths** (\textbackslash{}(D/\textbackslash{}lambda\textbackslash{})) sets the highest spherical-harmonic modes the optics can pass. Small \textbackslash{}(f\textbackslash{}\#\textbackslash{}) (large \textbackslash{}(D\textbackslash{})) → large \textbackslash{}(\textbackslash{}ell\_\{\textbackslash{}max\}\textbackslash{}) → finer angular details survive.

\paragraph{Canonical Equation.}
\[
D/\lambda
\]

\paragraph{Dictionary.}
Bessel + Aperture → Spectral Cutoff, Link to Spherical Harmonics: D/λ — D/λ

\paragraph{Encyclopedia.}
Bessel + Aperture → Spectral Cutoff, Link to Spherical Harmonics: D/λ is a symbolic expression, written D/λ. It introduces a symbol or relation used elsewhere in the framework. It is built from λ, D. Within the Trawin Zero Point registry it serves as one element of the framework's mathematical narrative.
